\section{Functions, closures, options, pointers}
\label{sec:functions}

\subsection{Functions}

A \ic{def} declares a function; parameters are typed, the return type is
inferred or checked. Write recursion with \ic{self} --- that is the
portable spelling, and it is what a function \emph{literal} (which has no
name) has to use. A \emph{named} \ic{def}'s own name is also in scope
inside its body, so \ic{count(n - 1, acc + 1)} and \ic{self(n - 1, ...)}
are the same call:

\begin{lstlisting}
def fact(n: long) -> long {
    if (n <= 1) { 1 } else { n * self(n - 1) }
};
\end{lstlisting}

\noindent Functions are first-class values: they may be bound
(\ic{let f = g;}), stored in arrays and records, passed to other
functions, and called through any of those (\ic{f(1)}, \ic{fs[0](1)},
\ic{r.handler(1)}). Calling a non-function is a compile error.

Function types include both parameter and return types. A function may also
be passed where the same parameters and a \ic{void} result are expected: the
caller discards its result. Other result types must agree.

\subsection{Tail calls}

A self-call in \emph{tail position} --- the body's trailing expression,
a \ic{return}'s value, or an if-arm's trailing expression --- recycles
the stack frame instead of growing it: the arguments are evaluated, the
frame is torn down, and control re-enters the function. Tail recursion
therefore runs in constant stack space:

\begin{lstlisting}
def sum_to(n: long, acc: long) -> long {
    if (n == 0) { acc } else { self(n - 1, acc + n) }
};
\end{lstlisting}

\noindent Calls outside tail position (\ic{self(x) + self(y)},
\ic{let t = self(x)}, inside loops) are ordinary calls and stack as
usual.

\subsection{Closures}
\label{sec:closures}

A function literal captures the enclosing bindings it reads --- and only
those that are already in scope where the literal is written (a name
bound by a \ic{let} \emph{after} the \ic{def} is not in scope and is an
undefined-variable error, exactly like reading a local before its
declaration). Captured bindings are \textbf{shared}: a write inside the
closure is visible outside, and a write outside is visible inside.

A literal written as an \emph{arrow} is the same closure with less
punctuation: \ic{v => v + 1} is \ic{def(v) \{ v + 1 \}},
\ic{(e, i) => e + i} takes two parameters, and \ic{() => 0} takes none.
The body is one expression, or a block; a parameter may carry a type
(\ic{(acc: long, v: long) => acc + v}), and the parentheses are optional
for a lone parameter. Arrow literals appear anywhere a \ic{def} literal
does, and they are what the array methods take
(Section~\ref{sec:array-methods}).

\begin{lstlisting}
def counter() {
    let n = 0;
    def tick() -> long { n = n + 1; n };
    tick
};
let t = counter();
t(); t();
print(t());        // 3 -- the one shared n
\end{lstlisting}

\noindent Each \emph{invocation} of the enclosing function creates fresh
bindings, so closures created in different iterations see different
cells (``\ic{let} semantics''):

\begin{lstlisting}
while (i < 3) {
    let x = i;
    fns += def() -> long { x };
    i = i + 1;
};
\end{lstlisting}

\noindent The \emph{loop variable} of a \ic{for x in a..b} or a
\ic{foreach (x in a)} is fresh on each iteration, like the \ic{let}
inside the loop above --- it is the one binding the loop makes per
iteration rather than one the body does:

\begin{lstlisting}
let fns: [() -> long] = [];
for i in 0..3 {
    fns += def() -> long { i };     // this iteration's i: 0, 1, 2
};
foreach (e in [7, 8, 9]) {
    fns += def() -> long { e };     // 7, 8, 9
};
\end{lstlisting}

\noindent A binding that is \emph{not} a loop variable --- one declared
before the loop, or the \ic{let} inside its body --- is a single cell,
and every closure the loop makes shares it, so they all see where the
loop ended:

\begin{lstlisting}
let k = 0;
while (k < 3) {
    k = k + 1;
    fns += def() -> long { k };     // one k, shared: every call says 3
};
\end{lstlisting}

\subsection{Options}
\label{sec:builtin-option}

\ic{T?} is the type of ``a \ty{T} or nothing''; its empty case is
\ic{none}. \ic{some(v)} boxes a value as present, \ic{is\_some(v)}
tests, and \ic{unwrap(v)} --- or its postfix spelling \ic{v!} ---
returns the payload and is a \emph{runtime error} on \ic{none} (it is a
programmer error, loudly). \ic{x == none} is the is-empty check, and
two options are equal when both are \ic{none} or both hold equal payloads (Section~\ref{sec:equality});
an option of an array is an optional reference, not a copied one.

\begin{lstlisting}
let line = readline();          // str?
if (line == none) { print("eof"); } else { print(line! + "!"); };
\end{lstlisting}

\subsection{Void}
\label{sec:void}

A \ic{def} whose result is not a value is \ic{void}. Write \ic{-> void}
to say so: the body's trailing value is then discarded, and \ic{return}
takes no value. An unannotated \ic{def} whose body ends in a statement
infers \ic{void}, and then every \ic{return} in it must be bare: a
\ic{return 5;} there is an error, and so is a bare \ic{return;} in a body
that ends in a value. A void call is a statement; using it as a value --- binding
it, printing it, passing it as an argument --- is a compile error:

\begin{lstlisting}
def log(msg: str) -> void { print("log: " + msg); };
log("hi");                 // fine: a statement
let x = log("hi");         // error: cannot bind a void value to 'x'
\end{lstlisting}

\subsection{Pointers}
\label{sec:pointers}

\ic{*T} is a raw address of one \ty{T}; \ic{*p = v} stores,
\ic{*p} reads. The address-of operator \ic{\&e} is deliberately narrow:

\begin{itemize}[nosep]
\item \ic{\&x} --- address of a local or parameter;
\item \ic{\&a[i]} --- address of an array element (any array expression:
  \ic{\&grid[0][i]} is legal);
\item \ic{\&arr} on a whole array --- \emph{both} kinds, \ic{[T]} and
  \ic{[T; N]} --- is a compile error
  (``cannot take the address of a whole array``), and \ic{\&} of any
  temporary is a different one (``'\&' needs a plain variable``). A binding names
  the array \emph{header}, not its first element, so \ic{\&arr} would hand
  out the header; write \ic{\&arr[0]}, which works for either kind.
\end{itemize}

\noindent Pointers are for interfacing with the system-call level
(buffers for \ic{read}/\ic{write}-style calls); ordinary code uses
arrays and records. They support C-style arithmetic: \ic{p + n} and
\ic{p - n} move by whole \ty{T} elements (one byte for \ic{*byte}, the
width of a sized integer for \ic{*u32} and its kind, eight bytes for every
other pointer), \ic{n + p} is the mirror,
\ic{p - q} counts elements when both pointers have the same type, and
\ic{+=}, \ic{-=}, \ic{++} and \ic{--} scale the same way. Ordering
comparisons compare raw addresses; equality is word equality. The
offset must be \ic{long} or \ic{byte}, and mixing pointer types in a
difference is a compile error.

\noindent\textbf{Sized integers.} To read or write memory laid out by
someone else (a system-call structure, a file or wire format), a pointer's
target can be a sized integer: \ic{i8}, \ic{i16}, \ic{u16}, \ic{i32},
\ic{u32} and \ic{u64}, and the 16-, 32- and 64-bit ones with a \ic{be}
suffix for big-endian (\ic{*u32be}); \ic{le} spells the native order and
is the same type. (\ic{u8} is \ty{byte} and \ic{i64} is \ty{long}.) A sized
integer is \emph{only} a pointer's target: \ic{*p} is a \ty{long}, sign- or
zero-extended from the width in memory, and \ic{*p = v} stores the low
bytes of \ic{v} and drops the rest, as a store to a \ic{byte} does. A
variable, parameter, field or array element of type \ic{u32} is a compile
error (``u32 is only a pointer's target (write *u32)''): the width is a
property of the memory, not of the value. The pointer comes from
\ic{ptr\_from}, below, or from any \ic{*T} of the same type.

\noindent\textbf{Packed records.} A record normally keeps each field in a
word of its own, sorted by name; \ic{packed} lays the fields out in the order
written, byte after byte with no padding, so a record can describe a region of
memory someone else defined:
\begin{verbatim}
type Flags = bits u8 { ready, error, mode: 3, _: 2 };
type Hdr = packed { magic: u32be, len: u16, flags: Flags, tag: [byte; 4] };
let h: Hdr = { magic: 0xCAFEBABE, len: 513, flags: { ready: true, mode: 5 } };
h.len = h.len + 1;
h.flags.error = true;
let p: *Hdr = ptr_from(addr);   // fields over memory you did not make
\end{verbatim}
\noindent Inside \ic{packed} the sized integers of the last paragraph are field
types (\ic{u32be}, \ic{i16}, \ic{u8}\ldots), as are \ic{long}, \ic{float},
\ic{bool} (one byte), a pointer, a fixed array \ic{[T; N]} of those, and another
packed record, which sits inline. Anything with a header of its own (\ic{str},
\ic{[T]}, a map, an ordinary record) is refused when the type is used: the
collector does not read raw memory, so a reference stored in it would not be
kept alive. \ic{bits u8 \{ a, b, mode: 3 \}} (backing \ic{u8}, \ic{u16},
\ic{u32} or \ic{u64}) gives each name the next bits from the least significant
end: one bit, read as a \ty{bool}, unless a width follows; \ic{\_} skips bits.
Writing one bit field reads and rewrites the whole backing integer.

\noindent A packed value \emph{is} the address of its bytes. \ic{h.len} reads
or writes the field in place (widened by its sign, byte-swapped for \ic{be});
\ic{h.flags} and \ic{h.sub} are views of the nested bytes, so they take no
assignment as a whole, only their fields; an array field reads as a \ic{*T} to its
first element. A literal \ic{\{ .. \}} where a packed type is expected allocates
a zeroed block of \ic{sizeof(Hdr)} bytes and stores the fields it names.
\ic{Hdr} and \ic{*Hdr} convert to each other, and \ic{ptr\_addr(h)} is the
address. \ic{sizeof(T)} and \ic{offsetof(T, field)} are the layout as constants.
A packed value is compared and copied as an address, and cannot be printed.

\noindent\textbf{Pointers and integers.} A pointer is its address at run
time, and two builtins say so. \ic{ptr\_addr(p)} is the \ic{long} value of a
pointer. \ic{ptr\_from(n)} is the pointer whose address is the integer
\ic{n} (a \ic{long} or a \ic{byte}), the inverse. Its result has no pointee
type of its own, so the type comes from where it is used: write
\ic{let p: *byte = ptr\_from(a)}, pass it to a \ic{*T} parameter, or store it
in a \ic{*T} field. A bare \ic{let p = ptr\_from(a)} is a compile error
(``ptr\_from does not say which pointer type it makes''), and so is a
non-integer argument. Nothing checks that \ic{n} is mapped memory: a wrong
address faults when it is dereferenced, which is what a kernel's physical
memory access needs and why ordinary code does not use it.

\noindent\textbf{A pointer does not outlive what it points into.}
\ic{\&x} on a local or a parameter is the address of a cell in that
function's frame, and the frame is reclaimed the moment the function
returns: a pointer to a local that escapes reads whatever the stack
holds by then --- a number that was never written, or an address, and
it changes from run to run. Returning one is not a type error,
because the checker cannot see how long a value lives; the rule is that
a pointer into a frame may only be used while that frame is live, which
means in practice that it must not be returned and must not be stored
in anything outliving the call. Taking the address of a \emph{global}
is fine (its storage is permanent), and so is \ic{\&a[i]}: a heap
array's elements and a fixed array's live in a heap cell, never in a
frame. To hand data to a function or keep it, return the value, an
array or a record --- those never dangle --- and take a pointer to
\ic{\&a[i]} where the callee needs an address.
